\subsection{Approximation spaces}

Given locally convex spaces E and F, recall (EVT, III, p. 14, example 2) that one denotes \(\mathcal{L}_{\mathrm{pc}}(\mathrm{E};\mathrm{F})\) the locally convex space obtained by equipping \(\mathcal{L}(\mathrm{E};\mathrm{F})\) with the topology of precompact convergence. We also denote \(\mathcal{L}_{\mathrm{pc}}(\mathrm{E})\) the locally convex space \(\mathcal{L}_{\mathrm{pc}}(\mathrm{E};\mathrm{E})\) . When the locally convex space E is Hausdorff and quasi-complete (EVT, III, p. 8, def. 6), the topology of precompact convergence on \(\mathcal{L}(\mathrm{E};\mathrm{F})\) coincides with the topology of compact convergence, since the closure of a precompact subset of E is compact (EVT, III, p. 8).

DEFINITION 2. --- A locally convex space E is said to have the approximation property, or again to be an approximation space, if the identity map \(1_{E}\) of E is adherent in the space \(\mathcal{L}_{\mathrm{pc}}(\mathrm{E})\) to the set \(\mathcal{L}^{\mathrm{f}}(\mathrm{E})\) of continuous finite-rank endomorphisms of E.

In particular, every finite-dimensional locally convex space is an approximation space.

Remarks. --- 1) For the topological direct sum of locally convex spaces E and F to be an approximation space, it is necessary and sufficient that E and F be approximation spaces.

\begin{enumerate}
\def\labelenumi{\arabic{enumi})}
\setcounter{enumi}{1}
\item
  Let E be a locally convex space. Let N be the closure of \(\{0\}\) in E and P an algebraic supplementary subspace of N in E. Then P is a topological supplement of N in E, and is isomorphic to the locally convex space E/N. The space N is an approximation space. By Remark 1, for E to be an approximation space, it is necessary and sufficient that the Hausdorff locally convex space E/N be one.
\item
  Let E be a locally convex space, A an equicontinuous subset of \(\mathcal{L}^{\mathrm{f}}(\mathrm{E})\) and T a total subset of E. If \(1_{\mathrm{E}}\) is adherent to A for the topology of simple convergence on T, then \(1_{\mathrm{E}}\) is adherent to A in \(\mathcal{L}_{\mathrm{pc}}(\mathrm{E})\) (EVT, III, p. 16, prop. 4 and p. 17, prop. 5), and E is an approximation space.
\item
  Let E be a locally convex space over C. Denote by \(E_{0}\) the locally convex space over R underlying E. For E to be an approximation space, it is necessary and sufficient that \(E_{0}\) itself be one. Indeed, the condition is necessary; let us prove that it is sufficient. Suppose therefore that \(E_{0}\) is an approximation space. Let C be a precompact subset of E and U a convex balanced neighborhood of 0 in E. Set \(C' = C \cup iC\) . There exists a continuous finite-rank R-linear map u from \(E_{0}\) into \(E_{0}\) such that \(x - u(x)\) belongs to U for all \(x \in C'\) . Set \(v(x) = \frac{1}{2}(u(x) - iu(ix))\) for all x in E. One thus defines a continuous C-linear map of finite rank from E into E. For every \(x \in C\) , we have \(x \in C'\) , \(ix \in C'\) and \(x - v(x) = \frac{1}{2}(x - u(x)) - \frac{i}{2}(ix - u(ix))\) , so that \(x - v(x)\) belongs to U. This proves that E is an approximation space.
\item
  Let E be a locally convex space over R. For the complexified locally convex space \(\mathrm{E}_{(\mathbf{C})}\) to be an approximation space, it is necessary and sufficient that E itself be one. This follows from Remarks 1 and 4, since the underlying real locally convex space of \(\mathrm{E}_{(\mathbf{C})}\) is isomorphic to \(E \times E\) .
\end{enumerate}

PROPOSITION 10. --- Let E be a Hilbert space. Then E is an approximation space.

The set A of finite-rank orthogonal projectors on E is equicontinuous. Let \(n \geqslant 1\) be an integer and \((x_{1}, \ldots, x_{n})\) be a family of elements of E. Let V denote the vector subspace generated by the \(x_{i}\) and \(p_{V}\) the orthogonal projector with image V. We have \(p_{\mathrm{V}}(x_{i}) = x_{i}\) for \(1 \leqslant i \leqslant n\) . This implies that \(1_{E}\) is adherent to A for the topology of simple convergence, and it follows that E is an approximation space (Remark 3).

Other examples of approximation spaces will be given in No. 7 of III, p. 20.

Lemma 2. --- Let E be a locally convex space. Then E is an approximation space if and only if, for every precompact subset C of E and every neighborhood U of 0 in E, there exist an integer \(n \geqslant 0\) , elements \(e_{1}, \ldots, e_{n}\) of E, and continuous linear forms \(f_{1}, \ldots, f_{n}\) on E, such that \(x - \sum_{i=1}^{n} f_{i}(x)e_{i}\) belongs to U for all \(x \in C\) .

The definition of the topology of precompact convergence shows that the condition is sufficient. Conversely, suppose that E is an approximation space. Let P be a topological supplementary space in E of the closure of \(\{0\}\) (cf. remark 2). The set A of \(u \in \mathcal{L}^{\mathrm{f}}(\mathrm{E})\) such that \(u(\mathrm{E}) \subset \mathrm{P}\) is then dense in \(\mathcal{L}^{\mathrm{f}}(\mathrm{E})\) . Since P is separated, every element \(u\) of A is of the form \(x \mapsto \sum_{i=1}^{n} f_i(x)e_i\) , where \(n \in \mathbf{N}\) , \(e_1, \ldots, e_n\) are elements of E and \(f_1, \ldots, f_n\) are continuous linear forms on E (EVT, I, p. 14, th. 2). This proves that the condition is necessary.

Remark 6. --- There exist Banach spaces which do not have the approximation property, as was demonstrated by P. ENFLO (A counterexample to the approximation problem in Banach spaces, Acta Math. 130 (1973), p. 309--317; cf. Exercise 25 of III, p. 112). This answered a question of S. BANACH (Théorie des opérations linéaires, Monografie Matematyczne I, Warszawa, 1932, remarks in Chapter VI, §1). See also the remark on p. 21.

PROPOSITION 11. --- Let E and F be locally convex spaces. Suppose that E is an approximation space.

\begin{enumerate}
\def\labelenumi{\alph{enumi})}
\item
  The set \(\mathcal{L}^{\mathrm{f}}(\mathrm{E};\mathrm{F})\) is dense in \(\mathcal{L}_{\mathrm{pc}}(\mathrm{E};\mathrm{F})\) ;
\item
  The set \(\mathcal{L}^{\mathrm{f}}(\mathrm{F};\mathrm{E})\) is dense in \(\mathcal{L}_{\mathrm{pc}}(\mathrm{F};\mathrm{E})\) ;
\item
  Let \(\mathfrak{S}\) be a set of bounded subsets of F and \(u \in \mathcal{L}(\mathrm{F};\mathrm{E})\) . Suppose that the image under \(u\) of every subset of F belonging to \(\mathfrak{S}\) is precompact. Then \(u\) is adherent to \(\mathcal{L}^{\mathrm{f}}(\mathrm{F};\mathrm{E})\) for the \(\mathfrak{S}\) -topology (EVT, III, p. 13).
\end{enumerate}

Let \(v\) be an element of \(\mathcal{L}(\mathrm{E};\mathrm{F})\) . The map \(\varphi\colon w\mapsto v\circ w\) from \(\mathcal{L}_{\mathrm{pc}}(\mathrm{E})\) into \(\mathcal{L}_{\mathrm{pc}}(\mathrm{E};\mathrm{F})\) is continuous (TG, X, p. 5, prop. 3). We have \(\varphi(1_{\mathrm{E}})=v\) and \(\varphi(\mathcal{L}^{\mathrm{f}}(\mathrm{E}))\subset\mathcal{L}^{\mathrm{f}}(\mathrm{E};\mathrm{F})\) , hence \(v\) is adherent to \(\mathcal{L}^{\mathrm{f}}(\mathrm{E};\mathrm{F})\) in \(\mathcal{L}_{\mathrm{pc}}(\mathrm{E};\mathrm{F})\) . This proves \(a\) ).

Similarly, under the hypotheses of \(c\) ), the map \(\psi: w \mapsto w \circ u\) from \(\mathcal{L}_{\mathrm{pc}}(\mathrm{E})\) into \(\mathcal{L}_{\mathfrak{S}}(\mathrm{F};\mathrm{E})\) is continuous (loc. cit.). We have \(\psi(1_{\mathrm{E}}) = u\) and \(\psi(\mathcal{L}^{\mathrm{f}}(\mathrm{E})) \subset \mathcal{L}^{\mathrm{f}}(\mathrm{F};\mathrm{E})\) , hence \(u\) is adherent to \(\mathcal{L}^{\mathrm{f}}(\mathrm{F};\mathrm{E})\) in \(\mathcal{L}_{\mathfrak{S}}(\mathrm{F};\mathrm{E})\) .

The image of a precompact subset of F under a continuous linear map from F into E is precompact, so b) follows from c).

COROLLARY. --- Let E be a separated approximation space and F a locally convex space. Every compact linear map from F into E is adherent to \(\mathcal{L}^{\mathrm{f}}(\mathrm{F};\mathrm{E})\) in the space \(\mathcal{L}(\mathrm{F};\mathrm{E})\) endowed with the topology of bounded convergence.

This follows from prop. 11, c), since the image of a bounded subset of F under a compact linear map from F into E is relatively compact in E (remark 1 of III, p. 2), hence precompact.

PROPOSITION 12. --- Let E be a locally convex space, I a set, and for each \(i \in I\) let \(F_{i}\) be a locally convex space and \(v_{i}: E \to F_{i}\) a continuous linear map with dense image. Suppose that for every neighborhood U of 0 in E, there exist \(i \in I\) and a neighborhood V of 0 in \(F_{i}\) such that \(v_{i}^{-1}(V) \subset U\) . If the \(F_{i}\) are approximation spaces, then E is one as well.

Let A be a precompact subset of E and U a neighborhood of 0 in E. By hypothesis there exist \(i \in I\) and a continuous seminorm \(p\) on \(F_i\) such that U contains \((p \circ v_i)^{-1}([0,1])\) . Set \(F = F_i\) , \(v = v_i\) and \(B = v(A)\) , and suppose that F is an approximation space. The set B is precompact in F. Hence, by Lemma 2, there exist an integer \(n \geqslant 1\) , elements \(y_1, \ldots, y_n\) of F and continuous linear forms \(f_1, \ldots, f_n\) on F, such that one has, for all \(y \in B\) ,

\[
p \left(y - \sum_ {j = 1} ^ {n} f _ {j} (y) y _ {j}\right) \leqslant \frac {1}{2}.
\]

Since B is bounded (EVT, III, p. 3, prop. 2), there exists a real number M > 0 such that \(|f_{j}(y)| \leqslant M\) for all j such that \(1 \leqslant j \leqslant n\) and every \(y \in B\) . Moreover, since \(v(\mathrm{E})\) is dense in F, there exists, for each integer j such that \(1 \leqslant j \leqslant n\) , an element \(x_{j}\) of E such that \(p(y_{j} - v(x_{j})) \leqslant (2n\mathrm{M})^{-1}\) . The linear map

\[
u \colon x \longmapsto \sum_ {j = 1} ^ {n} f _ {j} (v (x)) x _ {j}
\]

belongs to \(\mathcal{L}^{\mathrm{f}}(\mathrm{E})\) . For every \(x\in \mathrm{A}\) , we have

\[
v (x - u (x)) = v (x) - \sum_ {j = 1} ^ {n} f _ {j} (v (x)) y _ {j} + \sum_ {j = 1} ^ {n} f _ {j} (v (x)) \left(y _ {j} - v \left(x _ {j}\right)\right),
\]

whence \(p(v(x - u(x))) \leqslant \frac{1}{2} + \frac{nM}{2nM} = 1\) , and consequently \(x - u(x) \in U\) . The proposition follows.

COROLLARY 1. --- A dense vector subspace of an approximation space is an approximation space.

COROLLARY 2. --- If the separated completion of a locally convex space E is an approximation space, then E is an approximation space.

COROLLARY 3. --- The product of a family of approximation spaces is an approximation space.

Indeed, let \((\mathrm{E}_{i})_{i\in\mathrm{I}}\) be a family of approximation spaces. For every finite subset J of I, set \(E_{J}=\prod_{i\in J}E_{i}\) and let us denote \(v_{J}\) the canonical map from \(E=\prod_{i\in I}E_{i}\) into \(E_{J}\) . The locally convex space \(E_{J}\) is an approximation space (Remark 1). The corollary then follows from Prop. 12 applied to the locally convex space E, to the family \((\mathrm{E}_{\mathrm{J}})\) of locally convex spaces and to the continuous linear maps \(v_{J}:E\to E_{J}\) .

COROLLARY 4. --- Let E be a locally convex space. If every continuous seminorm on E is majorized by a continuous pre-Hilbertian seminorm, then E is an approximation space.

Let \(\mathcal{P}\) be the set of continuous pre-Hilbertian seminorms on E. For each \(p \in \mathcal{P}\) , the hypothesis implies that the seminormed space \(\mathrm{E}_p\) obtained by equipping E with the seminorm \(p\) is an approximation space (Corollary 2 and Proposition 10), and the identity mapping from E into \(\mathrm{E}_p\) is continuous. The corollary therefore follows from Proposition 12.

Lemma 3. --- Let E be a metrizable locally convex space and \((x_{n})_{n\in\mathbf{N}}\) a sequence of elements of E converging to 0. There exists a sequence \((y_{n})_{n\in\mathbf{N}}\) of elements of E converging to 0 and a sequence \((\lambda_{n})_{n\in\mathbf{N}}\) of elements of the interval [0,1] converging to 0 such that we have \(x_{n}=\lambda_{n}y_{n}\) for all \(n\in N\) .

Since E is metrizable, there exists a fundamental system \((\mathrm{V}_{m})_{m\in\mathbf{N}}\) of balanced neighborhoods of 0 in E such that \(V_{0}=E\) and \(2^{m+1}V_{m+1}\subset V_{m}\) for all \(m\geqslant0\) . Since \((x_{n})\) converges to 0, there exists a strictly increasing sequence \((\mathrm{N}_{m})_{m\in\mathbf{N}}\) of integers such that \(N_{0}=0\) and such that, for all \(m\geqslant0\) , one has \(x_{n}\in V_{m}\) for \(n\geqslant N_{m}\) . For every integer \(n\geqslant0\) there exists a unique integer \(m\geqslant0\) such that \(N_{m}\leqslant n<N_{m+1}\) , and we then set \(\lambda_{n}=2^{-m}\) and \(y_{n}=2^{m}x_{n}\) . The sequence \((\lambda_{n})\) thus defined converges to 0. Moreover, since \(y_{n}\in V_{m}\) for \(n\geqslant N_{m+1}\) the sequence \((y_{n})\) converges to 0 in E. Finally, we have \(x_{n}=\lambda_{n}y_{n}\) for all n.

Recall (EVT, II, p. 28) that if E is a vector space and L is a convex balanced subset of E, we denote by \(E_{L}\) the vector subspace of E generated by L, equipped with the seminorm whose unit ball is L. When the set L contains no line, this seminorm is a norm.

Lemma 4. — Let E be a Fréchet space and A a compact subset of E. There exists a compact convex balanced subset L of E containing A such that the topologies induced on A by those of E and of \(E_{L}\) coincide.

The set A is contained in the closed convex balanced hull of the set of points of a sequence \((x_{n})_{n\in\mathbf{N}}\) of elements of E converging to 0 (EVT, IV, p. 24, cor. 1). Let \((y_{n})_{n\in\mathbf{N}}\) and \((\lambda_{n})_{n\in\mathbf{N}}\) be sequences satisfying the conclusions of Lemma 3. The closed convex balanced hull L of the points of the sequence \((y_{n})\) contains A and is a compact subset of E (EVT, III, p. 8). Consequently, \(E_{L}\) is a Banach space (EVT, III, p. 8, cor.). In this Banach space, the norm of \(x_{n}\) is bounded above by \(\lambda_{n}\) , hence the sequence \((x_{n})\) tends to 0. The closed convex balanced hull \(\widetilde{A}\) of the sequence \((x_{n})\) in \(E_{L}\) is a compact subset of \(E_{L}\) (EVT, III, p. 8). As L is a bounded subset of E, the canonical injection of \(E_{L}\) into E is continuous. The topologies induced on \(\widetilde{A}\) by those of \(E_{L}\) and E coincide, and \(\widetilde{A}\) is a compact, hence closed, subset of E. Since the set \(\widetilde{A}\) is convex, balanced, and contains the sequence \((x_{n})\) ; it contains A. This completes the proof.

THEOREM 4. --- Let E be a Fréchet space.

\begin{enumerate}
\def\labelenumi{\alph{enumi})}
\item
  Suppose that E is an approximation space. Then, for every semi-normed space F, the space \(\mathcal{L}^{\mathrm{f}}(\mathrm{F};\mathrm{E})\) is dense in \(\mathcal{L}^{\mathrm{c}}(\mathrm{F};\mathrm{E})\) for the topology of bounded convergence;
\item
  Conversely, suppose that for every Banach space F, every compact linear map from F into E is adherent to \(\mathcal{L}^{\mathrm{f}}(\mathrm{F};\mathrm{E})\) for the topology of bounded convergence. Then E is an approximation space.
\end{enumerate}

Let F be a semi-normed space. The closure of \(\mathcal{L}^{\mathrm{f}}(\mathrm{F};\mathrm{E})\) in \(\mathcal{L}(\mathrm{F};\mathrm{E})\) is contained in \(\mathcal{L}^{\mathrm{c}}(\mathrm{F};\mathrm{E})\) (prop. 1 and 2 of III, p. 4). It is equal to \(\mathcal{L}^{\mathrm{c}}(\mathrm{F};\mathrm{E})\) if E is an approximation space by the cor. of prop. 11. This proves the assertion \(a\) ).

Suppose that the hypothesis of \(b\) holds. Let \(\varepsilon > 0\) be a real number. Let A be a compact subset of E and \(p\) a continuous seminorm on E. Let L be a compact convex balanced subset of E such that A is a compact subset of the normed space \(\mathrm{E_L}\) (Lemma 4). The canonical injection \(j: \mathrm{E_L} \to \mathrm{E}\) is compact and \(\mathrm{E_L}\) is a Banach space (EVT, III, p. 8, cor.). Hence, by hypothesis, there exist an integer \(n \geqslant 1\) , elements \(e_1, \ldots, e_n\) of E, and continuous linear forms \(\ell_1, \ldots, \ell_n\) on \(\mathrm{E_L}\) , such that the map \(v\) of \(\mathrm{E_L}\) to E defined by \(v(x) = \sum_{i=1}^{n} \ell_i(x)e_i\) satisfies \(p(x - v(x)) \leqslant \frac{\varepsilon}{2}\) for \(x \in A\) . The image of \({}^{t}j: E' \to (E_{L})'\) is dense in \((E_{L})'\) for the weak topology (EVT, IV, p. 6, prop. 5). On \((E_{L})'\) the topology of compact convergence is compatible with the duality between \((E_{L})'\) and \(E_{L}\) (EVT, IV, p. 3, example). Hence \({}^{t}j(E')\) is dense in \((E_{L})'\) for the topology of compact convergence (EVT, IV, p. 1, prop. 1), and there exist continuous linear forms \(f_{1}, \ldots, f_{n}\) on E such that

\[
| \ell_ {i} (x) - f _ {i} (x) | p (e _ {i}) \leqslant \frac {\varepsilon}{2 n}
\]

for \(x \in \mathrm{A}\) and \(1 \leqslant i \leqslant n\) . The endomorphism \(u \colon x \mapsto \sum_{i=1}^{n} f_i(x)e_i\) of \(\mathrm{E}\) belongs to \(\mathcal{L}^{\mathrm{f}}(\mathrm{E})\) and for every \(x \in \mathrm{A}\) , we have

\[
p (x - u (x)) \leqslant p (x - v (x)) + p (v (x) - u (x)) \leqslant \frac {\varepsilon}{2} + n \times \frac {\varepsilon}{2 n} = \varepsilon .
\]

It follows from this that \(1_{E}\) is adherent to \(\mathcal{L}^{\mathrm{f}}(\mathrm{E})\) for the topology of compact convergence. Now this topology coincides with the topology of precompact convergence because E is complete. Therefore E is an approximation space.

